\section{Critical inputs and normalization}\label{sec:inputs}

This section specifies the consequences of Assumption~\ref{ass:cardy}
that we import from the companion manuscript~\cite{V}. In particular,
the exploration and normalization statements below are previously
proved inputs, not additional hypotheses and not conclusions of scalar
convergence alone without further argument. The remainder of this paper
proves the annular estimates and the amplitude from these inputs.

Throughout, $C,c$ denote positive constants whose values may change
between occurrences. A comparison $u\asymp v$ means that $u/v$ lies
between two fixed positive constants. Fixed-factor changes in arm
radii and fixed buffer widths are absorbed in these constants. We use
unit-intensity geometry except when explicitly introducing a mesh
parameter. Let $b_*$ be a sufficiently large fixed microscopic radius.

\subsection{Arms, localization, and local changes}
An arm in an annulus is a monochromatic cell path joining its two
boundary components. Alternating four arms have colors black, white,
black, white in cyclic order. Let $A(r,s)$ be their probability in the
square-annulus convention of~\cite{V}, for $b_*\le r<s$.

\begin{proposition}[Critical arm and localization inputs]\label{in:arms}
There are $g\in(1,2)$, $\alpha,\beta>0$, and constants $C,c>0$ such
that, for $b_*\le r<s<t$,
\begin{align}
 c(r/s)^g&\le A(r,s)\le C(r/s)^{1+\alpha},\label{in:power}\\
 C^{-1}A(r,s)A(s,t)&\le A(r,t)\le C A(r,s)A(s,t).
 \label{in:quasimult}
\end{align}
The one-arm probability is at most $C(r/s)^\beta$. The probability of
three alternating arms in a straight half-plane is at most
$C(r/s)^2$.

For a collar of diameter comparable to $s$, an event determined by the
Poisson points in a slightly enlarged collar makes the graph and all
vertex stars in the smaller collar local to those points, with mesh
smaller than any prescribed fixed fraction of $s$. Its failure
probability is at most $C\exp(-cs^2)$. The widths and constants can
depend on that fraction. Such events may be chosen rotation equivariant
about the center. Events with disjoint enlarged supports use independent
Poisson variables.
\end{proposition}

These are Lemmas~2.1 and~2.3 of~\cite{V}. The quantitative arm theory
builds on the box-crossing and annealed arm estimates for Voronoi
percolation; see~\cite{Tassion,Vanneuville}. A localization event will
be called a \emph{screen}. For example, requiring that every ball of
a sufficiently small fixed relative radius in a collar contain a
Poisson point is rotation equivariant and has the asserted exponential
bound, by a finite covering and the Poisson empty-ball formula.
Extra buffers give locality of both cells and vertex stars. Thus a
point inserted beyond the enlarged support cannot send its cell through
the screened collar.

We will also use the necessary-arm conclusions of Lemma~2.2
of~\cite{V}. A color change or a local insertion that changes a
crossing requires four alternating arms out of its interior support.
Near a straight side, these continue as three alternating half-plane
arms; beyond the scale of a corner, at least one arm continues. These
statements use the exact closed-cell crossing convention. They apply
to a mixed difference of two local changes as well: if the mixed
difference is nonzero, each local change affects the appropriate
crossing in at least one of the other local states. Screens keep the
other modification out of the arm probes used for this implication.

For later independence arguments it is useful to allow failed screens
inside an arm probe. Here is the convention and the estimate behind it.
Place screens at geometrically spaced radii. If the largest failed
radius is $v$, retain only the arm test starting at $C_0v$, with its
screens, where $C_0$ provides a fixed unused gap. This retained test
is independent of the failed screen. Quasi-multiplicativity and
\eqref{in:power} give, when $C_0v<s$,
\[
 A(C_0v,s)\le C(v/r)^g A(r,s).
\]
Consequently, summing the failed-screen terms over geometric radii
costs at most
\[
 C A(r,s)\sum_{v\ge r}(v/r)^g e^{-cv^2}
 \le C A(r,s).
\]
If $C_0v\ge s$, the exponential bound and the polynomial lower bound
give the same conclusion. The corresponding construction for one-arm
and straight half-plane three-arm tests preserves their displayed
power upper bounds. We call these \emph{screened arm probes}. Their
supports can be separated by fixed-factor gaps, and they can be used
in the presence of inserted sites outside their buffers.

\subsection{Fixed-scale convergence and exploration domains}
The following inputs are qualitative. Their order of limits will matter:
all radii, buffer widths, and domain tests are fixed before the mesh
tends to zero.

\begin{proposition}[Scaling-limit inputs]\label{in:limits}
Under Assumption~\ref{ass:cardy}, the following consequences hold.
\begin{enumerate}[label=\textup{(\roman*)}]
\item From every sequence of Poisson meshes tending to zero one can
extract, after coupling, a further subsequence along which the graphs
almost surely form a deterministic good sequence in the sense of
\cite[Lemma~2.5]{V}. Along such a sequence, critical crossings of
Jordan quadrilaterals converge, also for graph domains whose boundaries
and marked points converge in curve order. Finite joint crossing tests
and tests with finitely many separated prescribed squares have the
common triangular-percolation limit.
\item Chordal color explorations in converging Jordan graph domains
converge to chordal $\mathrm{SLE}_6$. At a pre-query stop on a
crosscut, the remaining target face, its boundary banks, and its
conformal maps converge in the sense of the stopped-domain results
of~\cite{V}. A bank is an individual boundary arc in prime-end order;
distinct banks may have touching physical images. Each relevant bank
is individually simple.
\item For some $c_0>0$, in the mesh-first limit the probability of
four bulk arms with any specified color order is at most
$C(r/s)^{1+c_0}$, and that of six bulk arms with any specified color
order is at most $C(r/s)^{2+c_0}$. Three free arms at a regular
boundary, measured in a conformal chart, have an upper bound with
exponent strictly larger than one.
\end{enumerate}
\end{proposition}

For precise locators, part~(i) uses Proposition~2.4, Lemma~2.5, and
Proposition~4.5 of~\cite{V}; part~(ii) uses Proposition~3.1 and
Lemmas~3.4--3.7; part~(iii) uses Lemmas~2.7, 2.8, and~2.11.
The free boundary-arm estimate in Lemma~2.11 has exponent $4/3$.
``Free'' means that the arms use the independent fair interior sites,
not the sites in the prescribed base arc.

We emphasize two domain qualifications. First, a boundary chart must
converge uniformly on its closure, have no constant boundary subarc,
and project injectively on the compact bank interval where the
estimate is used. Separate prime-end banks need not have disjoint
physical projections. The local estimate applies in an annulus by
choosing a small graph disk patch with one side on that bank and the
remaining sides in the interior. These are Jordan approximations;
reflection in the parameter coordinate gives bounded distortion on
smaller patches. Second, product colors are conditional on the graph.
After an exploration, its remaining colors are fresh, but a
remaining-domain theorem can be used only after its geometric
hypotheses have been established. We never infer uniform convergence
after conditioning on arbitrary atypical explored colors. The uniform
operator estimate in Section~\ref{sec:transfer} will instead use
independent core variables and bounds on the signed weights.

We use the site triangulation of the Voronoi tiling from~\cite{V}.
Its faces are triangles, drawn through the tiling rather than required
to have straight edges. An interface strand, or \emph{wall}, passes
through a triangle between its two heterochromatic edges. Inside a
screened collar, simple site cycles approximate circles in curve
order. They may be constructed by following narrow strips and
retaining simple cuts; same-boundary chords on the domain-facing side
can be shortcut while retaining the specified inner region. Auxiliary
uniform angular choices, if needed, preserve rotation invariance.

All transfer cuts in this paper lie in the bulk. Their incident cells
and stars stay away from the quad boundary, where site connectivity
and cell connectivity could otherwise require additional care.
At the actual quad boundary, we retain the original closed-cell event
and use the necessary-arm estimates just stated. Null nongeneric
geometries, and null sets of insertion locations in integrals, are
discarded throughout.

\subsection{Calibration and the preliminary exponent}
Let $S_\eps$ be the probability that a site inserted at the center of
$Q_0=[-1,1]\times[-1/2,1/2]$ is color-pivotal for its left-to-right
crossing in intensity $\eps^{-2}$. Equivalently, $S_\eps$ is the
expectation of the black-forced crossing indicator minus the
white-forced crossing indicator for that site. Set
\[
 m_\eps=\eps^{-2},\qquad
 q_\eps=\bigl(\E N_\eps(Q_0)\bigr)^{-1},\qquad
 f(R)=S_{1/R}.
\]
The last equality agrees by Euclidean scaling with the unit-intensity
definition in the introduction.

\begin{proposition}[Calibration]\label{in:calibration}
There are constants $c,B\in(0,\infty)$ such that
\begin{equation}\label{in:calibration-eq}
 m_\eps q_\eps S_\eps\longrightarrow c,
 \qquad q_\eps\E N_\eps\longrightarrow B.
\end{equation}
Moreover,
\begin{equation}\label{in:comparison}
 f(R)\asymp A(b_*,R).
\end{equation}
\end{proposition}

The limits are precisely the single-time Poisson specialization of
Lemmas~7.1 and~7.2 of~\cite{V}, called ``Common calibration'' and
``The prescribed normalization.'' The signed color weight at the
inserted site is $\delta_1-\delta_0$; crossing monotonicity identifies
its expectation with $S_\eps$. The comparison follows from the local
pivotal bounds and the square-to-site calibration in Sections~2 and~6
of~\cite{V}. These inputs do not assert that
$R^{5/4}f(R)$ converges.

\begin{lemma}[Exponent without an amplitude]\label{in:exponent}
Under Assumption~\ref{ass:cardy},
\begin{equation}\label{in:exponent-eq}
 f(R)=R^{-5/4+o(1)}\qquad(R\to\infty).
\end{equation}
\end{lemma}
\begin{proof}
Force all cells meeting a central square of radius $h$ in $Q_0$
first black and then white, and denote the difference of crossing
probabilities by $D_\eps(h)$. Proposition~4.5 of~\cite{V} gives,
for each fixed small $h>0$, a common limit of this quantity in the
Voronoi and triangular models. The upper arm bound and
\cite[Lemma~6.1]{V}, ``A square can transmit a pivotal change,''
compare $D_\eps(h)$ with the four-arm probability from radius $h$
to radius $1$, with constants uniform in small $h$, after taking
the mesh limit. The corresponding triangular four-arm probabilities
have exponent $\lambda=5/4$ by~\cite{SW}.

Fix $\eta>0$. For an arbitrarily large but fixed ratio $L=1/h$,
these comparisons and qualitative convergence imply, for all
sufficiently large $r$ depending on $L$,
\[
 c_\eta L^{-\lambda-\eta/3}
 \le A(r,Lr)\le C_\eta L^{-\lambda+\eta/3}.
\]
Choose $L$ so large that both these constants and the constants in
\eqref{in:quasimult} are absorbed by an additional factor
$L^{\eta/3}$. Only after fixing this $L$ choose a starting radius
above the convergence threshold. Iterating over radii $r,Lr,L^2r,
\ldots$ gives four-arm bounds with exponents between
$\lambda-\eta$ and $\lambda+\eta$. The bounded initial range and
the final incomplete ratio change only the multiplicative constants,
by \eqref{in:power} and \eqref{in:quasimult}. Equation~\eqref{in:comparison}
then yields
$C_\eta^{-1}R^{-\lambda-\eta}\le f(R)\le C_\eta R^{-\lambda+\eta}$.
Letting $\eta\downarrow0$ proves the lemma.
\end{proof}

We will henceforth write $\lambda=5/4$. The remaining task is to
replace the $o(1)$ in this exponent by convergence of the multiplicative
coefficient. The next two analytic sections identify which annular
symmetries make that improvement possible.
